\begin{topic}[id=memfault_remote_debug,
             difficulty={Mittel},
             type={Tool- und Prozessanalyse},
             prerequisites={Grundlagen eingebetteter Systeme; Basiswissen Debugging und Software-Wartung},
             practice={gut möglich},
             owner={Dozententeam},
             updated={2026-04-09},
             tags={ Observability, Crashdump, Core-Dump, Telemetrie, OTA, Datenschutz }]{Remote-Debugging: Memfault}

\TopicDescription{Feldfehler sind teuer und schwer reproduzierbar. Memfault bietet Embedded-Observability: Geräte liefern Crash/Core-Dumps, Heartbeats und Telemetrie, die in der Cloud ausgewertet werden—inklusive OTA-Funktionen. Das Thema beschreibt Architektur und Integrationsschritte (Build/CI, Symbole/Backtraces) und zeigt, wie Ursachen sichtbar werden (Out-of-Memory, Stack-Überlauf, Watchdog). Ebenso wichtig: Datenschutz—welche Daten sind nötig, wie schützt man sie, wie vermeidet man PII? Ergebnis: Ein klares Verständnis, wann Observability den Unterschied macht.}

\TopicLeitfragen{\item Wie integriert man Symbolik/Backtraces sinnvoll in CI/CD?
\item Welche Telemetrie ist „mindestens nötig“ vs. „nice to have“?
\item Wie adressiert man Datenschutz/PII praktisch?
}
\TopicScope{\item Kein Vendor-Lock-in-Vergleich im Detail.
\item Keine vollständige OTA-Implementierung.
}
\TopicPractice{Simulation eines Crashes und Interpretation des Backtraces in der Plattform.}

\TopicKeywords{Observability, Crashdump, Core-Dump, Telemetrie, OTA, Datenschutz}

\nocite{memfaultDocs,memfaultBlogIntro}
\end{topic}
